\section{System Model, Threats and Invariants}
\label{sec:model}

\subsection{Parties}
A \emph{saga} is one unit of agent work: a plan of steps, each a named action on
a registered \emph{participant} (an agent, a tool, a validator), with a declared
\emph{effect class}: \code{PURE}, \code{REVERSIBLE}, \code{COMPENSABLE} (has a
registered compensating action), \code{IRREVERSIBLE\_GATED} (an effect the
outbox can hold) or \code{IRREVERSIBLE\_IMMEDIATE} (in the world as soon as the
participant returns). One \emph{writer} process, the daemon
\code{janus-orchd}, owns an evidence log and hosts the saga engine; agents,
consoles and interception proxies are its clients. \emph{Answerers} (validator
processes and people) decide the gates the engine cannot decide itself. An
\emph{auditor} holds a copy of the log (or an exported bundle) and the writer's
public key, and nothing else.

\subsection{Threats in scope}
\begin{itemize}
\item \textbf{After-the-fact tampering} with any record by anyone without the
  writer's signing key, including deletion, insertion and reordering. Cutting
  records off the \emph{end} of a bundle is detected when its manifest is signed,
  which an auditor can require (\code{-require-signed-manifest}), or against a
  head obtained out of band; an unsigned bundle verified without either reads as
  a shorter log.
\item \textbf{A gated agent authoring its own permission}: stating the approval,
  balance or recommendation a gate decides on, moving an approval it received
  onto a different proposal (\cref{sec:findings}), or recording an answer, a
  declaration or a result under another participant's name or a person's.
\item \textbf{Prompt injection} that makes a model \emph{want} an out-of-policy
  effect. Janus does not try to stop the wanting; it bounds what the wanting can
  release.
\item \textbf{A crash anywhere}: the writer, the coordinator or a participant
  killed between any two records, including between an approval and the decision
  it permits.
\end{itemize}

\subsection{Out of scope, stated}
\label{sec:model-out}
\begin{itemize}
\item \textbf{Participants without keys, and calls that record nothing.} A
  participant whose manifest declares no key is taken at its word unless the
  daemon requires every caller to sign; replication, operator, effect-delivery
  and read-only calls are not authenticated, and the channel is not encrypted.
\item \textbf{What a signature does not cover.} A declaration's signature covers
  its saga, step and facts but not its attempt, so a captured declaration can be
  replayed into a later attempt of the same step. An answer's signature covers
  its verdict, not the facts the answerer was shown: the fold binds the answer to
  the pinned proposal, so an answerer deceived about the facts over the
  unencrypted channel approves the real proposal. A key withdrawn from a manifest
  is still accepted for sagas that pinned the version declaring it, until they
  end.
\item \textbf{A writer-key holder writing a consistent false history.} Segment
  signatures prove who wrote a log, not that it is the only log that writer
  produced. External anchoring of log heads is designed but not built; today an
  auditor can pin an expected head obtained out of band (\code{-expect-head}).
\item \textbf{Semantic fidelity of a declaration.} A gate decides on the amount a
  step \emph{declares}. If a model extracts 4{,}000 from a request for 40{,}000,
  every gate here judges 4{,}000 (\cref{sec:eval-model}).
\item \textbf{Byzantine replicas} and multi-writer consensus: a follower that
  diverges is detected and attributed; a lease fences a superseded writer; neither
  tolerates a malicious one.
\item \textbf{Clients that ignore a refusal} on the SDK path (\cref{sec:design-edges}).
\end{itemize}

\subsection{Invariants}
\Cref{tab:invariants} lists the eight invariants and where each is enforced.
They are checked by property tests, by crash-injection suites that kill the
writer at every transition, and by an adversarial ``evil auditor'' suite that
performs each tampering attack for real.

\begin{table}[t]
\caption{Invariants and where they are enforced}
\label{tab:invariants}
\centering\footnotesize
\begin{tabular}{@{}p{0.9cm}p{3.1cm}p{3.6cm}@{}}
\toprule
& Invariant & Enforced by \\
\midrule
I1 & Evidence before effect: no release without a durable chained append & the outbox writes \code{RELEASING} durably before delivering \\
I2 & Chain integrity: any byte mutation is detected & segment format; the offline verifier \\
I3 & Compensation closure: every non-\code{PURE} step has a compensation or a gate & gate admission \\
I4 & Commit safety: a gated irreversible effect releases only from \code{COMMITTED}, at least once, under one idempotency key & outbox and idempotency keys \\
I5 & Replay determinism: transitions are a pure function of the event sequence & the saga fold \\
I6 & Frontier safety: no commit over an unsealed conflicting touch & the frontier gate, where the policy has one \\
I7 & Compensation order: reverse-topological & the compensation planner \\
I8 & Registry pinning: every participant version resolves at replay & registry admission \\
\bottomrule
\end{tabular}
% source: pkg/outbox, pkg/evidence, pkg/gate (admission), pkg/saga (fold, compensation planner, frontier), pkg/registry
\end{table}
